\section{BGP Graceful Restart, LLGR, and Route Withdrawal Behaviour}

\subsection{Objective}

The objective of this experiment was to test BGP control-plane restart behaviour and route withdrawal behaviour.

This experiment has two parts:

\begin{enumerate}
\item BGP Graceful Restart / LLGR forwarding continuity test
\item BGP route withdrawal test
\end{enumerate}

In the Graceful Restart test, the BGP session was restarted, but the route was expected to be retained temporarily so that forwarding could continue.

In the withdrawal test, the route itself was intentionally removed from BGP export. In that case, forwarding should eventually fail because the route is no longer valid.

\subsection{BGP GR and LLGR Configuration Verification}

Before running the test, the BGP configuration was checked on the BGP routers.

The following BGP options were already configured:

\begin{verbatim}
graceful restart yes;
long lived graceful restart yes;
\end{verbatim}

The protocol details also showed:

\begin{verbatim}
Long-lived graceful restart
LL stale time: 3600
Restart time: 120
BGP state: Established
\end{verbatim}

The BGP sessions were established before testing.

\begin{table}[h]
\centering
\begin{tabular}{|l|l|l|}
\hline
\textbf{Router} & \textbf{BGP Peer} & \textbf{State} \\
\hline
\texttt{hpe-r1} & \texttt{hpe-r9} & Established \\
\hline
\texttt{hpe-r1} & \texttt{hpe-r2} & Established \\
\hline
\texttt{hpe-r2} & \texttt{hpe-r9} & Established \\
\hline
\texttt{hpe-r9} & \texttt{hpe-r1} & Established \\
\hline
\texttt{hpe-r9} & \texttt{hpe-r2} & Established \\
\hline
\end{tabular}
\caption{BGP session state before GR and withdrawal testing}
\end{table}

\subsection{BGP Graceful Restart Forwarding Test}

\subsubsection{Need for Isolation}

A simple BGP restart test is not enough to prove Graceful Restart. If the \texttt{hpe-r1 -> hpe-r9} BGP path fails, traffic might simply move through another path:

\begin{verbatim}
hpe-r1 -> hpe-r2 -> hpe-r9
\end{verbatim}

That would prove BGP failover, not Graceful Restart. Therefore, the final test used isolation.

\subsubsection{Isolation Method}

The final Graceful Restart proof used a blackhole guard and disabled the alternate path.

The setup was:

\begin{verbatim}
Traffic: hpe-h1 -> hpe-h3
Source host: 10.0.61.2
Destination host: 10.0.93.2
Target prefix: 10.0.93.0/24
Expected BGP next-hop: 10.0.19.3
\end{verbatim}

A blackhole guard route was added on \texttt{hpe-r1}:

\begin{verbatim}
blackhole 10.0.93.0/24 metric 9999
\end{verbatim}

At the same time, the valid BGP route still existed:

\begin{verbatim}
10.0.93.0/24 via 10.0.19.3 dev eth0 proto bird metric 32
\end{verbatim}

The Linux route lookup still preferred the real BGP route:

\begin{verbatim}
10.0.93.2 via 10.0.19.3 dev eth0
\end{verbatim}

Then the alternate \texttt{hpe-r1 -> hpe-r2} path was disabled:

\begin{verbatim}
hpe-r1 eth1 DOWN
hpe-r1 protocol r2 disabled/down
\end{verbatim}

This means the test had no usable alternate path through \texttt{hpe-r2}.

\subsubsection{BGP Restart Action}

The BGP session from \texttt{hpe-r9} toward \texttt{hpe-r1} was restarted:

\begin{verbatim}
docker exec hpe-r9 birdc restart r1
\end{verbatim}

During the test, the BGP session on \texttt{hpe-r1} temporarily became non-established:

\begin{verbatim}
BGP non-established seen ms: 114
State: Active
Received: Administrative reset
\end{verbatim}

This proves that a real BGP control-plane restart happened.

\subsubsection{GR Forwarding Result}

\begin{table}[h]
\centering
\begin{tabular}{|l|r|}
\hline
\textbf{Measurement} & \textbf{Result} \\
\hline
Blackhole guard enabled & yes \\
\hline
\texttt{hpe-r1 -> hpe-r2} link down & yes \\
\hline
Old next-hop after restart & yes \\
\hline
Alternate next-hop after restart & no \\
\hline
Blackhole/unreachable after restart & no \\
\hline
BGP non-established seen & 114 ms \\
\hline
Ping transmitted & 200 \\
\hline
Ping received & 200 \\
\hline
Ping lost & 0 \\
\hline
Packet loss & 0\% \\
\hline
\end{tabular}
\caption{BGP GR forwarding continuity result}
\end{table}

The important result lines were:

\begin{verbatim}
Old next-hop seen after restart: yes
Alternate next-hop seen after restart: no
Blackhole/unreachable seen after restart: no
Ping transmitted: 200
Ping received: 200
Ping lost: 0
Ping loss percent: 0
\end{verbatim}

\subsubsection{GR Interpretation}

This confirms that forwarding continued during the BGP restart.

The result is strong because the \texttt{hpe-r1 -> hpe-r2} alternate path was disabled, and a blackhole guard was installed for \texttt{10.0.93.0/24}. The traffic did not fall into the blackhole route. Instead, route lookup continued to use the old direct next-hop \texttt{10.0.19.3}, and no ping packets were lost.

Therefore, the traffic continuity was not caused by alternate routing or static/default fallback. It was caused by the forwarding plane continuing to use the retained route during the BGP restart.

The sampled BIRD route output did not show a visible `stale'' or `LLGR stale'' label. Therefore, the conclusion is not that BIRD visibly marked the route as stale. The correct conclusion is that forwarding continuity was observed during BGP restart under a GR/LLGR-enabled configuration.

\subsection{BGP Route Withdrawal Test}

\subsubsection{Purpose}

The withdrawal test proves the opposite behaviour.

In this test, the BGP session stayed up, but the prefix \texttt{10.0.93.0/24} was intentionally withdrawn from export toward \texttt{hpe-r1}.

The purpose was to show that Graceful Restart protects forwarding during restart, but it does not keep forwarding alive when the route is truly withdrawn.

\subsubsection{Withdrawal Isolation Setup}

The same isolation idea was used:

\begin{verbatim}
blackhole guard enabled
hpe-r1 -> hpe-r2 alternate path disabled
\end{verbatim}

Before withdrawal, \texttt{hpe-r1} still had the valid BGP route:

\begin{verbatim}
10.0.93.0/24 via 10.0.19.3 dev eth0 proto bird metric 32
blackhole 10.0.93.0/24 metric 9999
\end{verbatim}

Route lookup still preferred the real BGP route:

\begin{verbatim}
10.0.93.2 via 10.0.19.3 dev eth0
\end{verbatim}

The alternate path was disabled:

\begin{verbatim}
hpe-r1 eth1 DOWN
\end{verbatim}

A baseline ping before withdrawal succeeded:

\begin{verbatim}
5 packets transmitted
5 received
0\% packet loss
\end{verbatim}

\subsubsection{Withdrawal Action}

The \texttt{hpe-r9} BGP export policy toward \texttt{hpe-r1} was changed to reject the host-side prefix:

\begin{verbatim}
10.0.93.0/24
\end{verbatim}

The modified configuration was applied using:

\begin{verbatim}
docker exec hpe-r9 birdc configure
\end{verbatim}

\subsubsection{Withdrawal Result}

\begin{table}[h]
\centering
\begin{tabular}{|l|r|}
\hline
\textbf{Measurement} & \textbf{Result} \\
\hline
Blackhole guard enabled & yes \\
\hline
\texttt{hpe-r1 -> hpe-r2} path disabled & yes \\
\hline
Withdrawn prefix & \texttt{10.0.93.0/24} \\
\hline
BGP session non-established after withdrawal & no \\
\hline
Alternate next-hop after withdrawal & no \\
\hline
BIRD route missing time & 134 ms \\
\hline
Ping transmitted & 220 \\
\hline
Ping received & 11 \\
\hline
Ping lost & 209 \\
\hline
Packet loss & 95\% \\
\hline
Final ping after withdrawal & 100\% packet loss \\
\hline
\end{tabular}
\caption{BGP route withdrawal result}
\end{table}

The important result lines were:

\begin{verbatim}
BGP non-established seen after withdrawal: no
BIRD route missing ms: 134
BIRD route missing line: Network not found
Ping transmitted: 220
Ping received: 11
Ping lost: 209
Ping loss percent: 95
\end{verbatim}

Final ping after withdrawal:

\begin{verbatim}
5 packets transmitted
0 received
100\% packet loss
\end{verbatim}

\subsubsection{Withdrawal Interpretation}

This proves that the route was actually withdrawn while the BGP session stayed up.

The session did not go down:

\begin{verbatim}
BGP non-established seen after withdrawal: no
\end{verbatim}

But the route disappeared from BIRD:

\begin{verbatim}
Network not found
\end{verbatim}

After the route disappeared, traffic failed.

This is the expected behaviour. Graceful Restart helps during a control-plane restart, but it should not keep forwarding traffic forever for a route that has been intentionally withdrawn.

\subsection{Comparison}

\begin{table}[h]
\centering
\scriptsize
\renewcommand{\arraystretch}{1.25}
\begin{tabular}{|p{3cm}|p{3cm}|p{4.6cm}|p{3.2cm}|}
\hline
\textbf{Test} & \textbf{BGP Session} & \textbf{Route Behaviour} & \textbf{Traffic Result} \\
\hline
GR forwarding test & Temporarily Active & Forwarding continued via old next-hop & 0\% loss \\
\hline
Route withdrawal test & Stayed established & Route became \texttt{Network not found} & 95\% loss during test \\
\hline
\end{tabular}
\caption{Comparison of BGP Graceful Restart and route withdrawal}
\end{table}

\subsection{Conclusion}

The BGP Graceful Restart / LLGR experiment showed that forwarding continued during a BGP control-plane restart. During the restart, the BGP session temporarily moved into a non-established state, but traffic continued through the old direct next-hop \texttt{10.0.19.3} with 0\% packet loss.

The test was isolated using a blackhole guard and by disabling the alternate \texttt{hpe-r1 -> hpe-r2} path. Therefore, the result was not caused by alternate routing or default-route fallback.

The route withdrawal test showed the opposite behaviour. When \texttt{10.0.93.0/24} was intentionally withdrawn from BGP export, the BGP session stayed up, but the route disappeared from BIRD within 134 ms. Traffic then failed, with 95\% packet loss during the test and 100\% packet loss in the final ping.

Together, these two tests show the difference between Graceful Restart route retention and true route withdrawal.
